
\subsubsection{Proposed Protocol}

The experimental design specifies:

\begin{itemize}
\item \textbf{Dataset}: LeanDojo Benchmark (98,734 theorems from Lean mathlib v3.4.2)
\item \textbf{Model}: LeanDojo ReProver (ByT5-based retrieval-augmented transformer)
\item \textbf{Sample}: 100 theorems randomly sampled from test set
\item \textbf{Timeout}: Extended timeout of 300 seconds (100× the standard 3 seconds)
\item \textbf{Confidence Window}: First 15 proof steps
\item \textbf{Ground Truth}: Theorems proven within 300s labeled "successful"; those timing out labeled "timeout"
\end{itemize}

\subsubsection{Evaluation Metrics}

\begin{itemize}
\item Pearson correlation (r) and Spearman rank correlation (ρ)
\item Area Under ROC Curve (AUC)
\item F1 Score, Precision, Recall for detector variants
\item Statistical significance testing (α=0.05)
\end{itemize}

\subsubsection{Actual Implementation}

As documented in the research directory, the experimental code includes:

\begin{enumerate}
\item Mock data generation functions in \texttt{test_mock.py} that create synthetic results with predetermined statistical properties
\item A results file \texttt{results_raw.csv} containing 100 entries labeled "mock_theorem_XXX"
\item Variance values generated from normal distributions with specified means and standard deviations
\item LeanDojo integration code with placeholders but no evidence of actual execution with real theorems
\end{enumerate}

The mock data generator creates variance distributions matching the target hypothesis:
\begin{itemize}
\item Success group: mean=0.09, std=0.05
\item Timeout group: mean=0.29, std=0.18
\end{itemize}

These predetermined distributions ensure that subsequent statistical analysis will produce strong correlations and group differences.
